\documentclass[12pt]{article}
\usepackage{amsmath,amssymb,amsthm}
\usepackage{fullpage}

\newtheorem{theorem}{Theorem}[section]
\newtheorem{proposition}[theorem]{Proposition}
\newtheorem{lemma}[theorem]{Lemma}
\newtheorem{corollary}[theorem]{Corollary}
\newtheorem{definition}[theorem]{Definition}
\newtheorem{convention}[theorem]{Convention}
\newtheorem{remark}[theorem]{Remark}

\newcommand{\R}{\mathbb{R}}
\newcommand{\Z}{\mathbb{Z}}
\newcommand{\eps}{\varepsilon}
\newcommand{\Sig}{\Sigma}
\DeclareMathOperator{\interior}{int}

\title{Piecewise-Linear M\"obius Band Embeddings in $\R^3$:\\
the threshold is $\sqrt3$}
\author{}
\date{}

\begin{document}
\maketitle

\noindent
We prove that
\[
\inf\{\beta>0:\ \beta\text{ is realized}\}=\sqrt3 .
\]
The proof has two halves.
\begin{enumerate}
\item[\textbf{(S1)}] If $\beta$ is realized then $\beta>\sqrt3$.
\item[\textbf{(S2)}] For every $\eps>0$ there is a realized $\beta<\sqrt3+\eps$.
\end{enumerate}
Together these give that the infimum equals $\sqrt3$ (and that it is not
attained). The argument follows the strategy of R.~Schwartz, \emph{The
Optimal Paper M\"obius Band} (Annals of Math.\ 2025); the present write-up
is self-contained.

%=====================================================================
\section{Setup, definitions, and the reading of condition (iii)}
\label{sec:setup}

\paragraph{The strip, the group, the quotient.}
Let $\Sig=\R\times[0,1]$ with the identification
$(x,0)\sim(x+1,1)$, and let
\[
G=G_\beta:\Sig\to\Sig,\qquad G(x,y)=(x+\beta,\,1-y).
\]
$G$ is a Euclidean glide reflection: it is the composition of the
reflection $y\mapsto 1-y$ in the centreline $\R\times\{1/2\}$ with the
translation by $(\beta,0)$. It is an isometry of the flat metric on
$\Sig$, it has no fixed points (a fixed point would need $x=x+\beta$),
and the cyclic group $\langle G\rangle\cong\Z$ acts freely and properly
discontinuously: $G^k(x,y)=(x+k\beta,y)$ for even $k$ and
$(x+k\beta,1-y)$ for odd $k$, so any compact set meets only finitely
many of its translates. Hence the quotient $\Sig/G$ is a smooth flat
surface; since $G$ reverses the transverse orientation of the strip,
$\Sig/G$ is non-orientable, i.e.\ a flat M\"obius band. Its core curve is
the image of the centreline, which has length $\beta$. (The original
identification $(x,0)\sim(x+1,1)$ only records that $\Sig$ is itself a
flat M\"obius band of core length $1$; it plays no further role, since
$G_\beta$ already determines the quotient we study.)

A map $F:\Sig\to\R^3$ that is \emph{invariant} under $G$
(condition (i)), i.e.\ $F\circ G=F$, descends to a continuous map
$\overline F:\Sig/G\to\R^3$. Condition (ii), \emph{injectivity on
orbits} ($F(p)=F(q)\Rightarrow q\in G\cdot p$), says exactly that
$\overline F$ is injective. Thus (i)+(ii) say that $\overline F$ is a
topological embedding of the flat M\"obius band of core length $\beta$
into $\R^3$.

\paragraph{Triangulations adapted to the boundary.}
$F$ is piecewise linear, so it is affine on each triangle of some
locally finite triangulation $\mathcal T$ of $\Sig$. By passing to a
$G$-invariant common refinement we may and do assume:
\begin{itemize}
\item $\mathcal T$ is $G$-invariant (so it descends to a finite
triangulation of $\Sig/G$);
\item each vertex of $\mathcal T$ lying on $\partial\Sig=\R\times\{0,1\}$
is treated as a $\mathcal T$-vertex, and each triangle has its closure
either disjoint from $\partial\Sig$ or meeting $\partial\Sig$ in a single
edge.
\end{itemize}
We call such a $\mathcal T$ \emph{boundary-adapted}.

We will not need the full triangulation; we need only the following two
geometric notions, which we now define \emph{separately} (these were
conflated in earlier drafts).

\begin{definition}[Clean triangle; distinguished data; crosscuts]
\label{def:clean}
A triangle $\tau$ of a boundary-adapted triangulation is \emph{clean} if
exactly one of its edges lies in $\partial\Sig$. For a clean triangle:
\begin{itemize}
\item the \emph{distinguished edge} of $\tau$ is its edge in
$\partial\Sig$;
\item the \emph{distinguished vertex} of $\tau$ is the opposite vertex
(which lies on the \emph{other} boundary component of $\Sig$, because
$\tau$ has its base on one boundary line and its apex must reach across
or up; in our constructions and in the relevant region this is the
case --- see Remark \ref{rem:clean});
\item the two edges other than the distinguished edge are the
\emph{crosscutting edges};
\item a \emph{crosscut} of $\tau$ is a straight segment from the
distinguished vertex $p$ of $\tau$ to a point $q$ of the (relative
interior of the) distinguished edge. Equivalently, if $a,b$ are the
endpoints of the distinguished edge, a crosscut is
$\{(1-s)p+sq:\ 0\le s\le1\}$ with $q=(1-r)a+rb$, $r\in(0,1)$.
\end{itemize}
A boundary-adapted triangulation all of whose triangles are clean is
called a \emph{clean triangulation}.
\end{definition}

\begin{convention}[Reading of the squeezing condition (iii)]
\label{conv:squeeze}
A \emph{boundary subarc of a triangle} $\tau$ is a non-degenerate
subsegment of $\partial\tau$, i.e.\ a subsegment of one of the three
edges of $\tau$. Among the three edges of a clean triangle, only the
distinguished edge lies in the boundary $\partial\Sig$ of the M\"obius
band; the two crosscutting edges are interior chords shared with
neighbouring triangles. The squeezing condition (iii) ---
``\emph{every non-degenerate boundary subarc of a triangle is strictly
shortened by $F$}'' --- is the requirement that the affine map
$F|_\tau$ be \emph{strictly length-decreasing along each edge of} $\tau$
\emph{that lies in} $\partial\Sig$.

We adopt the standard convention, used throughout the paper-M\"obius-band
literature and forced by the construction in (S2), that a realizing map
is a \emph{squeeze map}: on each clean triangle $\tau$, $F|_\tau$
\emph{strictly shortens the distinguished edge and strictly lengthens
each crosscutting edge}. We make two remarks justifying that this is the
correct and intended reading.

\smallskip
\emph{(a) The boundary part.} Strict shortening of the distinguished edge
is exactly condition (iii) under the literal reading above (the only
edge in $\partial\Sig$ is the distinguished one).

\smallskip
\emph{(b) The crosscut part.} The crosscutting edges are not in
$\partial\Sig$, so (iii) imposes nothing on them directly. We require, in
addition, that they be lengthened. This is the defining feature of a
``squeeze'' (a developable/paper-folding type map: the band is pressed
flat, so chords transverse to the core are stretched while the boundary
is gathered/shortened). Both the lower bound (S1) and the construction
(S2) are stated and proved for squeeze maps in this sense, so the value
$\sqrt3$ is the threshold for the class of squeeze maps. We flag
explicitly that property (b) is part of the hypothesis on $F$, not a
consequence of (a) alone; indeed (a) by itself does not pin down a
threshold (one may shrink everything and embed a tiny band).
\end{convention}

\begin{remark}\label{rem:clean}
In all triangles used below (both in S1, where we only use the abstract
structure, and in S2, where every triangle is built explicitly), the
distinguished vertex of a clean triangle lies on the boundary component
opposite to its distinguished edge. This is automatic for the explicit
triangles of (S2) and is the only configuration we invoke.
\end{remark}

%=====================================================================
\section{Proof of Statement 1: $\beta>\sqrt3$}
\label{sec:S1}

Fix a realizing $F$ for $\beta$ with a $G$-invariant clean triangulation
$\mathcal C$, and write $f=F$. We freely use Convention
\ref{conv:squeeze}.

\paragraph{The crosscut foliation.}
The (relative) interiors of all crosscuts of all clean triangles foliate
$\interior(\Sig)$:

\begin{lemma}\label{lem:foliation}
Through every point $z\in\interior(\Sig)$ there passes exactly one
crosscut, and crosscuts in distinct triangles meet, if at all, only along
shared crosscutting edges.
\end{lemma}

\begin{proof}
$\interior(\Sig)$ is the disjoint union of the (open) triangle interiors,
the (open) crosscutting edges, and the $\mathcal C$-vertices in
$\interior(\Sig)$ (there are none on the open edges or in faces). It
suffices to handle each stratum.

\emph{Interior of a triangle.} Let $\tau$ be a clean triangle with apex
$p$ and base $[a,b]$. The radial projection from $p$ maps
$\tau\setminus\{p\}$ onto $[a,b]$, sending $z$ to the unique point
$q(z)$ where the ray $\overrightarrow{pz}$ meets $[a,b]$. For
$z\in\interior(\tau)$ we have $q(z)$ in the relative interior of $[a,b]$,
and the segment $[p,q(z)]$ is the unique crosscut of $\tau$ through $z$.
No crosscut of another triangle enters $\interior(\tau)$ because triangle
interiors are disjoint.

\emph{Open crosscutting edge $e$.} Let $e$ be shared by adjacent clean
triangles $\tau,\tau'$ (or, if $e\subset\partial\Sig$ it is a
distinguished edge and not an interior point; so assume $e$ interior).
The edge $e$ joins a base endpoint of $\tau$ to its apex $p$; thus $e$
is itself the crosscut of $\tau$ obtained by taking $q$ to be that base
endpoint --- a degenerate ($r\to0$ or $r\to1$) limit. More usefully, a
point $z$ in the open edge $e$ lies on the boundary crosscut of
\emph{both} $\tau$ and $\tau'$. We declare the crosscut through such a
$z$ to be $e$ itself (it is a genuine segment with endpoints on the two
boundary components, since one endpoint of $e$ is an apex on one boundary
line and the other endpoint of $e$ is a base vertex on the other). This
is consistent: $e$ is the common limiting crosscut of the two
neighbouring radial families, so the foliation closes up across $e$.

In all cases the crosscut through $z$ is unique, and two crosscuts in
different triangles share at most a common crosscutting edge. This is the
asserted (singular) foliation.
\end{proof}

\paragraph{Admissible pairs and the involution $I$.}
We parametrise crosscuts by a single real coordinate. For a crosscut $A$
let $m(A)$ be the $x$-coordinate of the point where $A$ crosses the
centreline $\R\times\{1/2\}$ (every crosscut meets the centreline
exactly once, since it runs from one boundary line to the other).
Write $A\prec B$ if $m(A)<m(B)$. Since $G$ acts on crosscuts (it is an
isometry preserving the triangulation, hence permuting clean triangles
and their crosscuts) and $G$ shifts the centreline coordinate by $\beta$,
we have $m(G^kA)=m(A)+k\beta$.

\begin{definition}
An ordered pair $(A,B)$ of crosscuts is \emph{admissible} if
$A\prec B\prec G(A)$, i.e.\ $m(A)<m(B)<m(A)+\beta$. Let $\mathcal A$ be
the set of admissible pairs.
\end{definition}

Admissibility forces $A$ and $B$ to lie in distinct $G$-orbits: if
$B=G^kA$ then $m(B)-m(A)=k\beta\in\{0,\beta,2\beta,\dots\}$ would have to
lie in the open interval $(0,\beta)$, impossible. Define
\[
I(A,B)=(B,\,G(A)).
\]
Then $(A,B)\in\mathcal A\iff I(A,B)\in\mathcal A$, because
$m(A)<m(B)<m(A)+\beta$ is equivalent to
$m(B)<m(GA)=m(A)+\beta<m(B)+\beta=m(GB)$. In the $(x,y)=(m(A),m(B))$
coordinates,
\[
\mathcal A=\{(x,y):x<y<x+\beta\},\qquad I(x,y)=(y,\,x+\beta),
\]
which is the affine glide reflection generating
$(x,y)\mapsto(x+\beta,y+\beta)$ as its square; in particular $I$ has no
fixed point in $\mathcal A$ and acts freely.

\begin{lemma}[Disjoint images]\label{lem:adm}
If $(A,B)$ is admissible then $f(A)$ and $f(B)$ have disjoint relative
interiors.
\end{lemma}

\begin{proof}
Suppose not: there are $p\in\interior(A)$, $q\in\interior(B)$ with
$f(p)=f(q)$. By condition (ii), $q=G^k(p)$ for some $k\in\Z$. Then
$q\in G^k(A)\cap B$ lies in $\interior(B)$ and in $G^k(\interior A)$,
which is the (relative) interior of the crosscut $G^k(A)$. By Lemma
\ref{lem:foliation} the crosscut through $q$ is unique, so $B=G^k(A)$.
Hence $m(B)-m(A)=k\beta\in(0,\beta)$, impossible for an integer $k$.
\end{proof}

\begin{lemma}[Squeeze maps expand crosscuts]\label{lem:expand}
$f$ strictly increases the length of every crosscut. More precisely, for
every clean triangle $\tau$ and every crosscut $A$ of $\tau$ one has
$\ell(f(A))>\ell(A)$.
\end{lemma}

\begin{proof}
\emph{Reduction to a single triangle.} By Definition \ref{def:clean}
every crosscut $A$ lies in a single clean triangle $\tau$ and runs from
the distinguished vertex (apex) $p$ of $\tau$ to a point $q$ of the
distinguished (base) edge. Thus the hypotheses needed below — that the two
crosscutting edges $[p,a],[p,b]$ are the lengthened edges and the
distinguished edge $[a,b]$ is the shortened one — apply verbatim, because
$A$ is the segment $[p,q]$ pointing from the apex toward the base. In
particular \emph{no assembly across several triangles is required for a
crosscut}: a crosscut never traverses more than one $\mathcal C$-triangle.
(The segments that do run through several triangles — the trapezoid lids
$C,D$ of Lemma \ref{lem:bound} — lie in $\partial\Sig$ and are treated
separately there, as \emph{shortened} boundary arcs, not as crosscuts.)
It therefore suffices to prove $\ell(f(A))>\ell(A)$ for $A=[p,q]\subset\tau$.

\smallskip
\emph{Linearisation.} Let $\phi=f|_\tau$, an affine map. Translate so that
$p=0\in\R^2$ and $f(p)=0\in\R^3$; then $\phi$ is linear. Put $V=a,\ W=b$
(as vectors from $p$, where $a,b$ are the endpoints of the distinguished
edge) and $V'=\phi(V),\ W'=\phi(W)$. The crosscutting edges $[p,a],[p,b]$
have lengths $\|V\|,\|W\|$, their images $[f(p),f(a)]=[0,V'],\,[0,W']$ have
lengths $\|V'\|,\|W'\|$, and the distinguished edge $[a,b]$ and its image
have lengths $\|V-W\|,\|V'-W'\|$. By Convention \ref{conv:squeeze} (the
crosscutting edges are strictly lengthened, the distinguished edge
strictly shortened),
\begin{equation}\label{eq:squeeze}
\|V\|<\|V'\|,\qquad \|W\|<\|W'\|,\qquad \|V-W\|>\|V'-W'\|.
\end{equation}
A crosscut of $\tau$ ending at $q=(1-r)a+rb$ with $r\in[0,1]$ is the
segment $\{s\,((1-r)V+rW):0\le s\le1\}$ from $p=0$ to
$(1-r)V+rW$, of length $\|(1-r)V+rW\|$; its image is the segment from
$0$ to $\phi((1-r)V+rW)=(1-r)V'+rW'$, of length $\|(1-r)V'+rW'\|$
(here we used linearity of $\phi$). So it suffices to prove
\begin{equation}\label{eq:goal}
\|(1-r)V+rW\|<\|(1-r)V'+rW'\|\qquad\text{for all }r\in(0,1).
\end{equation}
(The relevant range for a genuine crosscut is $r\in(0,1)$, i.e.\ $q$ in
the open base; the degenerate endpoints $r=0,1$ give the two crosscutting
\emph{edges} $[p,a],[p,b]$ themselves — these are the limiting crosscuts
appearing in the foliation of Lemma \ref{lem:foliation} — and for those
\eqref{eq:goal} reads $\|V\|<\|V'\|$ resp.\ $\|W\|<\|W'\|$, which are the
first two inequalities of \eqref{eq:squeeze}. Thus the strict expansion
in fact holds for the whole closed range $r\in[0,1]$, edges included.)

\smallskip
\emph{The key scalar inequality.} We first record
\begin{equation}\label{eq:dot}
V\!\cdot\!W<V'\!\cdot\!W'.
\end{equation}
By the polarisation identity $2\,V\!\cdot\!W=\|V\|^2+\|W\|^2-\|V-W\|^2$
(and likewise for the primed vectors), squaring the three inequalities of
\eqref{eq:squeeze} (all quantities are nonnegative, so squaring preserves
strict inequalities) gives $\|V\|^2<\|V'\|^2$, $\|W\|^2<\|W'\|^2$ and
$\|V-W\|^2>\|V'-W'\|^2$, whence
\[
2\,V\!\cdot\!W=\|V\|^2+\|W\|^2-\|V-W\|^2
<\|V'\|^2+\|W'\|^2-\|V'-W'\|^2=2\,V'\!\cdot\!W'.
\]
Each of the three summands moves strictly in the favourable direction
(the two squared lengths increase, the subtracted squared length
decreases), so the inequality is strict; this is \eqref{eq:dot}. Note all
three hypotheses of \eqref{eq:squeeze} are genuinely used here, and the
sign in front of $\|V-W\|^2$ is negative, which is exactly why the
hypothesis $\|V-W\|>\|V'-W'\|$ enters with the correct sign.

\smallskip
\emph{Term-by-term comparison.} Fix $r\in(0,1)$. Expanding the squared
norms,
\[
\|(1-r)V+rW\|^2=(1-r)^2\|V\|^2+r^2\|W\|^2+2r(1-r)\,(V\!\cdot\!W),
\]
\[
\|(1-r)V'+rW'\|^2=(1-r)^2\|V'\|^2+r^2\|W'\|^2+2r(1-r)\,(V'\!\cdot\!W').
\]
The three coefficients $(1-r)^2,\ r^2,\ 2r(1-r)$ are \emph{strictly
positive} for $r\in(0,1)$. Multiplying the strict inequalities
$\|V\|^2<\|V'\|^2$, $\|W\|^2<\|W'\|^2$ (from \eqref{eq:squeeze}) and
$V\!\cdot\!W<V'\!\cdot\!W'$ (from \eqref{eq:dot}) by these positive
coefficients and adding the three resulting strict inequalities yields
\[
\|(1-r)V+rW\|^2<\|(1-r)V'+rW'\|^2 .
\]
(No term has the wrong sign: every coefficient is positive and every
factor inequality points the same way, so the sum is strict.) Taking
nonnegative square roots gives \eqref{eq:goal}. Hence
$\ell(f(A))=\|(1-r)V'+rW'\|>\|(1-r)V+rW\|=\ell(A)$, as claimed.
\end{proof}

\paragraph{The topological lemma.}
For an oriented crosscut $C$ (oriented so as to point from the bottom
boundary toward the top boundary), let $\overrightarrow f(C)\in\R^3$ be
the vector along $f(C)$ with $\|\overrightarrow f(C)\|=\ell(f(C))$ and
pointing in the direction of the oriented image. Because $G$ reverses the
vertical orientation, $G(C)$ has the opposite vertical orientation, and
since $f\circ G=f$,
\begin{equation}\label{eq:odd-vec}
\overrightarrow f(G(C))=-\,\overrightarrow f(C),\qquad
f(\,\text{centreline pt of }G(C)\,)=f(\,\text{centreline pt of }C\,).
\end{equation}
Write $c(C)$ for the centreline point of $C$, so $f(c(GC))=f(c(C))$.

\begin{lemma}[T]\label{lem:T}
There is an admissible pair $(A,B)$ with $f(A)\perp f(B)$ (perpendicular
directions) and $f(A),f(B)$ coplanar.
\end{lemma}

\begin{proof}
On $\overline{\mathcal A}=\{x\le y\le x+\beta\}$ define
\[
g(A,B)=\overrightarrow f(A)\cdot\overrightarrow f(B),\qquad
h(A,B)=\big(f(c(A))-f(c(B))\big)\cdot\big(\overrightarrow f(A)\times
\overrightarrow f(B)\big).
\]
Both are continuous in $(m(A),m(B))$ because the crosscut, its image
direction, and its centreline image all vary continuously with the
centreline coordinate (the crosscut foliation of Lemma
\ref{lem:foliation} is continuous, and $f$ is continuous and PL).
Moreover:
\begin{itemize}
\item $g(A,B)=0$ iff the directions $f(A),f(B)$ are perpendicular;
\item $h(A,B)=0$ iff $f(A),f(B)$ are coplanar (the scalar triple product
of the join vector and the two directions vanishes).
\end{itemize}
So it suffices to show $(0,0)\in\Phi(\mathcal A)$ where $\Phi=(g,h)$.

\emph{Oddness $\Phi\circ I=-\Phi$.} Using \eqref{eq:odd-vec} with
$I(A,B)=(B,GA)$:
\[
g(B,GA)=\overrightarrow f(B)\cdot\overrightarrow f(GA)
=\overrightarrow f(B)\cdot(-\overrightarrow f(A))=-g(A,B),
\]
\[
h(B,GA)=\big(f(c(B))-f(c(GA))\big)\cdot\big(\overrightarrow f(B)\times
\overrightarrow f(GA)\big)
=\big(f(c(B))-f(c(A))\big)\cdot\big(\overrightarrow f(B)\times
(-\overrightarrow f(A))\big),
\]
and since $u\times(-v)=-(u\times v)=v\times u$ and the join vector flips
sign too,
$h(B,GA)=\big(f(c(A))-f(c(B))\big)\cdot\big(\overrightarrow f(A)\times
\overrightarrow f(B)\big)=-h(A,B)$. Thus $\Phi\circ I=-\Phi$.

\emph{Boundary behaviour.} On $\partial_0\mathcal A=\{y=x\}$ we have
$A=B$ (same crosscut), so the two directions coincide and the join vector
is $0$; hence $h=0$ and $g=\|\overrightarrow f(A)\|^2>0$. (Strict
positivity uses $\ell(f(A))>0$, which holds because by Lemma
\ref{lem:expand} $f$ lengthens crosscuts, so $f(A)$ is non-degenerate.)
On $\partial_1\mathcal A=\{y=x+\beta\}$ we have $B=G(A)$, so by
\eqref{eq:odd-vec} the directions are opposite: $g=-\|\overrightarrow
f(A)\|^2<0$, and again $h=0$ (here the join vector $f(c(A))-f(c(GA))=0$ by
\eqref{eq:odd-vec}, so $h$ vanishes identically). Therefore
\[
\Phi(\partial_0\mathcal A)\subset X_+:=\{(s,0):s>0\},\qquad
\Phi(\partial_1\mathcal A)\subset X_-:=\{(s,0):s<0\}.
\]

\emph{Winding argument.} For a continuous path
$p:[0,1]\to\R^2\setminus\{0\}$ let $\theta(p)\in\R$ denote its total
angular displacement about the origin (the change in a continuous branch
of the argument along $p$; this is well-defined and additive under
concatenation). Suppose, for contradiction, that
$(0,0)\notin\Phi(\mathcal A)$, so $\Phi(\mathcal A)\subset\R^2\setminus
\{0\}$; together with $\Phi(\partial_0\mathcal A)\subset X_+$ and
$\Phi(\partial_1\mathcal A)\subset X_-$ this means $\Phi$ maps
$\overline{\mathcal A}$ into $\R^2\setminus\{0\}$.

Let $\gamma:[0,1]\to\overline{\mathcal A}$ be any path with
$\gamma(0)\in\partial_0\mathcal A$, $\gamma(1)\in\partial_1\mathcal A$,
and $\gamma((0,1))\subset\mathcal A$. Then $P:=\Phi\circ\gamma$ is a path
in $\R^2\setminus\{0\}$ with $P(0)\in X_+$ (argument $\equiv 0$) and
$P(1)\in X_-$ (argument $\equiv\pi$). Hence its total angular
displacement satisfies $\theta(P)\in\pi+2\pi\Z$, i.e.
\begin{equation}\label{eq:half}
\tfrac{1}{2\pi}\,\theta(P)\in\tfrac12+\Z .
\end{equation}

Consider $\gamma^*:=I\circ\gamma$. Since $I$ preserves $\mathcal A$ and
maps $\partial_0\mathcal A\to\partial_1\mathcal A$ and
$\partial_1\mathcal A\to\partial_0\mathcal A$, the reversed path
$\overline{\gamma^*}(s):=\gamma^*(1-s)$ is again a path of the same type
(from $\partial_0\mathcal A$ to $\partial_1\mathcal A$ with interior in
$\mathcal A$). The strip $\overline{\mathcal A}$ is convex, so $\gamma$
and $\overline{\gamma^*}$ are homotopic \emph{through paths of this type}:
the straight-line homotopy
$H(s,r)=(1-r)\gamma(s)+r\,\overline{\gamma^*}(s)$ stays in
$\overline{\mathcal A}$, keeps $H(0,r)\in\partial_0\mathcal A$ and
$H(1,r)\in\partial_1\mathcal A$ (each boundary line is convex and both
endpoints of $\gamma,\overline{\gamma^*}$ lie on the same respective
line — after, if necessary, sliding the endpoints along the boundary
lines, which does not leave them), and has interior in $\mathcal A$ for
$r\in[0,1]$. Applying the continuous map $\Phi$ (valued in
$\R^2\setminus\{0\}$ by assumption) gives a homotopy in
$\R^2\setminus\{0\}$ between $\Phi\circ\gamma=P$ and
$\Phi\circ\overline{\gamma^*}$, through paths whose endpoints stay on
$X_+$ and $X_-$ respectively. Since the rays $X_\pm$ have fixed arguments
$0,\pi$, the quantity $\theta(\cdot)$ is invariant under such a homotopy;
hence
\[
\theta(P)=\theta(\Phi\circ\overline{\gamma^*}).
\]

Now use the oddness $\Phi\circ I=-\Phi$. We have
$\Phi\circ\gamma^*=\Phi\circ I\circ\gamma=-\,\Phi\circ\gamma=-P$, and
reversing reverses the sign of $\theta$, so
\[
\theta(\Phi\circ\overline{\gamma^*})
=-\,\theta(\Phi\circ\gamma^*)
=-\,\theta(-P)
=-\,\theta(P),
\]
where $\theta(-P)=\theta(P)$ because the central reflection $x\mapsto -x$
is a rotation by $\pi$, which preserves angular displacement. Combining,
$\theta(P)=-\theta(P)$, i.e.\ $\theta(P)=0$. This contradicts
\eqref{eq:half}. Therefore $(0,0)\in\Phi(\mathcal A)$, which yields an
admissible $(A,B)$ with $f(A)\perp f(B)$ and coplanar.
\end{proof}

\paragraph{The trapezoid estimate.}
Let $(A,B)$ be as in Lemma \ref{lem:T}. By Lemma \ref{lem:adm}, $f(A)$
and $f(B)$ have disjoint relative interiors; they are perpendicular and
coplanar. Hence either the line through $f(A)$ misses $\interior f(B)$ or
the line through $f(B)$ misses $\interior f(A)$. Replacing $(A,B)$ by
$I(A,B)$ if necessary (still admissible, still satisfying Lemma
\ref{lem:T}), we may assume the line through $f(A)$ is disjoint from
$\interior f(B)$.

Pre- and post-composing with isometries (which change none of the
relevant lengths) place the picture as follows:
\begin{itemize}
\item $f(A)$ lies on the $x$-axis of a plane $\R^2\subset\R^3$;
\item $f(B)$ lies on the non-positive part of the $y$-axis;
\item $f$ sends the top (apex) endpoint of $A$ to the left endpoint of
$f(A)$, and the top endpoint of $B$ to the top endpoint of $f(B)$ (which
is the origin or a point on the negative $y$-axis).
\end{itemize}
Let $t\in\R$ be defined by: the crosscut $A$, oriented upward in $\Sig$,
has slope $1/t$ (so $A$ rises by $1$ in $y$ while moving by $t$ in $x$;
$t$ is the signed horizontal displacement of $A$ across the unit-height
strip). Then
\[
\ell(A)=\sqrt{1+t^2}.
\]
Also $\ell(B)\ge1$, because $B$ is a crosscut and so spans the full unit
height of the strip, whence its length is at least $1$.

Let $C\subset\R\times\{1\}$ and $D\subset\R\times\{0\}$ be the boundary
segments joining the endpoints of $A$ to the corresponding endpoints of
$G(A)$. The four segments $A,\,C,\,G(A),\,D$ bound a trapezoid $T$; since
$G$ is a glide reflection, $T$ is an isosceles trapezoid and $G$
identifies opposite vertices of $T$. Consequently $f$ identifies those
opposite vertices as well (as $f\circ G=f$). We label the four vertices,
clockwise starting from the bottom-left, as
\[
a_2\ (\text{bottom-left}),\quad a_1\ (\text{top-left}),\quad
a_2'\ (\text{top-right}),\quad a_1'\ (\text{bottom-right}),
\]
with $G(a_k)=a_k'$ for $k=1,2$. (This is the labelling of the original
paper; the point is that $G$ identifies vertices in \emph{opposite}
corners of $T$, namely $a_2\!\leftrightarrow\!a_2'$ and
$a_1\!\leftrightarrow\!a_1'$, which sit at diagonally opposite corners
because $G$ is an orientation-reversing glide reflection.) Concretely the
left side is $A=[a_2,a_1]$, the right side is $G(A)=[a_2',a_1']$, the top
lid is $C=[a_1,a_2']$, and the bottom lid is $D=[a_2,a_1']$. Let $b_2$ be
the bottom endpoint of $B$.

Because $f\circ G=f$ and $G(a_2)=a_2'$, $G(a_1)=a_1'$, we have
$f(a_2')=f(a_2)$ and $f(a_1')=f(a_1)$. Hence:
\begin{itemize}
\item $f(A)=[f(a_2),f(a_1)]$ has endpoints $f(a_2),f(a_1)$;
\item $f(C)$ joins $f(a_1)$ to $f(a_2')=f(a_2)$ --- i.e.\ $f(C)$ is a path
between the \emph{two endpoints of} $f(A)$;
\item $f(D)$ joins $f(a_2)$ to $f(a_1')=f(a_1)$ --- i.e.\ $f(D)$ is also a
path between the two endpoints of $f(A)$.
\end{itemize}
(It is essential here that the lid $C$ runs from $a_1$ to $a_2'$ and the
lid $D$ from $a_2$ to $a_1'$; the diagonal identification
$a_k\!\leftrightarrow\!a_k'$ then sends \emph{each} lid to a path between
the two \emph{distinct} endpoints $f(a_1)\ne f(a_2)$ of $f(A)$. Had the
two top corners been $G$-related --- which they are not, since $G$
identifies diagonal corners --- the lids would have become null-homotopic
loops and the estimate below would fail.)

\begin{lemma}\label{lem:bound}
$\ell(C)>\sqrt{1+t^2}$ and $\ell(D)>\sqrt{5+t^2}$.
\end{lemma}

\begin{proof}
By Lemma \ref{lem:expand}, $f$ lengthens crosscuts; by Convention
\ref{conv:squeeze} (the boundary part / condition (iii)), $f$ shortens
boundary arcs. Hence
\[
\ell(f(A))>\ell(A)=\sqrt{1+t^2},\quad
\ell(f(B))>\ell(B)\ge1,\quad
\ell(C)>\ell(f(C)),\quad \ell(D)>\ell(f(D)).
\]
As established just above (using $f\circ G=f$ and the diagonal
identification $a_k\!\leftrightarrow\!a_k'$), $f(C)$ is a path joining the
two \emph{distinct} endpoints $f(a_2),f(a_1)$ of the straight segment
$f(A)$. A path joining the two endpoints of a straight segment has length
at least that of the segment, so $\ell(f(C))\ge\ell(f(A))$. Altogether
\[
\ell(C)>\ell(f(C))\ge\ell(f(A))>\ell(A)=\sqrt{1+t^2}.
\]

For $D$: the point $f(b_2)=(0,y)$ with $y\le0$ splits $f(D)$ into two
subarcs $\gamma_1,\gamma_2$. Reflect $\gamma_2$ in the horizontal line
$\{Y=y\}$ to get $\gamma_2^*$; reflection preserves length, and
$\gamma_1\cup\gamma_2^*$ is a path joining the same pair of points that
$f(D)$ joins horizontally, but now with both pieces on the same side, so
\[
\ell(f(D))=\ell(\gamma_1)+\ell(\gamma_2)
=\ell(\gamma_1)+\ell(\gamma_2^*)
\ge\sqrt{\,(\Delta x)^2+(\Delta Y)^2\,}.
\]
Here the endpoints of $f(D)$ are $f(a_2)$ and $f(a_1')=f(a_1)$, i.e.\
exactly the two endpoints of $f(A)$; since $f(A)$ lies on the $x$-axis,
both endpoints of $f(D)$ have $Y$-coordinate $0$. The reflection of
$\gamma_2$ across $\{Y=y\}$ moves the endpoint lying on the $x$-axis from
$Y=0$ to $Y=2y$, while leaving the other endpoint (on $\{Y=y\}$ side
shared with $\gamma_1$) and the point $(0,y)$ fixed. Hence
$\gamma_1\cup\gamma_2^*$ joins one endpoint of $f(A)$ (at $Y=0$) to the
reflection of the other (at $Y=2y$); these two points are separated
horizontally by exactly $\Delta x=\ell(f(A))>\sqrt{1+t^2}$ (the horizontal
extent of $f(A)$) and vertically by $|\Delta Y|=|2y|$. Finally, $f(b_2)$
is the bottom endpoint of $f(B)$, whose top endpoint lies at $Y\le0$ on
the negative $y$-axis; therefore $|y|=|Y\text{-coord of }f(b_2)|\ge
\ell(f(B))$, giving $|2y|\ge 2\,\ell(f(B))>2$. Therefore
\[
\ell(D)>\ell(f(D))\ge\sqrt{\,\ell(f(A))^2+(2\ell(f(B)))^2\,}
>\sqrt{(1+t^2)+4}=\sqrt{5+t^2}. \qedhere
\]
\end{proof}

\paragraph{Endgame.} We compute $\ell(C),\ell(D)$ exactly. Both are
straight segments, on the lines $y=1$ and $y=0$ respectively. Model the
crosscut $A$ as the segment from a bottom point $(x_b,0)$ to a top point
$(x_t,1)$, so its horizontal displacement is $t=x_t-x_b$ (slope $1/t$,
length $\sqrt{1+t^2}$). Since $G(x,y)=(x+\beta,1-y)$, the segment $G(A)$
runs from the top point $(x_b+\beta,1)$ to the bottom point
$(x_t+\beta,0)$. The boundary segment $C$ on $y=1$ joins $A$'s top point
to $G(A)$'s top point, and $D$ on $y=0$ joins $A$'s bottom point to
$G(A)$'s bottom point, so
\[
\ell(C)=\big|(x_b+\beta)-x_t\big|=|\beta-t|,\qquad
\ell(D)=\big|(x_t+\beta)-x_b\big|=|\beta+t|.
\]
Lemma \ref{lem:bound} therefore reads
\begin{equation}\label{eq:two}
|\beta-t|>\sqrt{1+t^2},\qquad |\beta+t|>\sqrt{5+t^2}.
\end{equation}

We first remove the absolute values. If $\beta\le t$ then
$|\beta-t|=t-\beta$, and \eqref{eq:two} gives
$t-\beta>\sqrt{1+t^2}>|t|\ge t$, hence $-\beta>0$, contradicting
$\beta>0$; so $\beta>t$. If $\beta\le -t$ then $|\beta+t|=-t-\beta$, and
\eqref{eq:two} gives $-t-\beta>\sqrt{5+t^2}>|t|\ge -t$, hence
$-\beta>0$, again a contradiction; so $\beta>-t$. Thus $\beta>|t|$, and
\eqref{eq:two} becomes
\[
\beta-t>\sqrt{1+t^2},\qquad \beta+t>\sqrt{5+t^2},
\]
i.e.
\[
\beta>u(t):=\sqrt{1+t^2}+t,\qquad \beta>v(t):=\sqrt{5+t^2}-t .
\]
Hence $\beta>\max\{u(t),v(t)\}$. Now $u$ is strictly increasing and $v$
strictly decreasing on $\R$ (indeed $u'(t)=\tfrac{t}{\sqrt{1+t^2}}+1>0$
because $|t|<\sqrt{1+t^2}$, and $v'(t)=\tfrac{t}{\sqrt{5+t^2}}-1<0$
because $|t|<\sqrt{5+t^2}$), and at $t_0=1/\sqrt3$,
\[
u(t_0)=\sqrt{1+\tfrac13}+\tfrac1{\sqrt3}=\tfrac{2}{\sqrt3}+\tfrac1{\sqrt3}
=\sqrt3,\qquad
v(t_0)=\sqrt{5+\tfrac13}-\tfrac1{\sqrt3}=\tfrac{4}{\sqrt3}-\tfrac1{\sqrt3}
=\sqrt3 .
\]
Because $u$ increases and $v$ decreases and both equal $\sqrt3$ at $t_0$,
we have $\max\{u(t),v(t)\}\ge\sqrt3$ for every $t\in\R$ (if $t\ge t_0$
then $u(t)\ge u(t_0)=\sqrt3$; if $t\le t_0$ then $v(t)\ge v(t_0)=\sqrt3$).
Therefore $\beta>\sqrt3$. This proves Statement 1. \hfill$\square$

%=====================================================================
\section{Proof of Statement 2: realizing $\beta<\sqrt3+\eps$}
\label{sec:S2}

Fix $\eps\in(0,1/100)$. We build an embedded polyhedral M\"obius band
$\Omega\subset\R^3$ from three equilateral triangles joined cyclically by
four thin connector rectangles, then read off a clean triangulated flat
M\"obius band $\Omega'=\Sig/G_\beta$ with $\beta<\sqrt3+\eps$ and an
embedding $\overline f:\Omega'\to\Omega$ satisfying conditions (i)--(iii).

Let $\mu=2/\sqrt3$; note the identity $\tfrac32\mu=\sqrt3$. ($\mu$ is the
side length of a ``clean equilateral'' triangle: an equilateral triangle
whose height is $1$, since the height of an equilateral triangle of side
$\mu$ is $\tfrac{\sqrt3}{2}\mu=1$.)

\paragraph{The range $\Omega$.}
Let $T_0\subset \R\times(-\infty,0]\times\{0\}$ be the equilateral
triangle of side $\mu+\eps^2$ with one edge centred at the origin on the
$x$-axis (the apex pointing in the $-y$ direction). Denote its left,
right, top edges by $\partial_-T_0,\partial_+T_0,\partial_0T_0$. Let
$E_0=\partial_0T_0+(0,\eps^3,0)$ (a small parallel push-off of the top
edge). Set $T_\pm=T_0\pm(0,0,\eps^3)$, two congruent copies in the planes
$z=\pm\eps^3$, with edges $\partial_jT_\pm$ parallel to $\partial_jT_0$
for $j\in\{0,+,-\}$.

Define three \emph{connectors}:
\begin{itemize}
\item $R_\pm=\operatorname{conv}(\partial_\pm T_0\cup\partial_\pm T_\pm)$,
a thin rectangle with long sides $\mu+\eps^2$ and short sides $\eps^3$;
\item $R_0=S_+\cup S_-$, where
$S_\pm=\operatorname{conv}(\partial_0T_\pm\cup E_0)$ are thin rectangles
with long sides $\mu+\eps^2$ and short sides $<2\eps^3$.
\end{itemize}
$R_\pm$ connects $\partial_\pm T_0$ to $\partial_\pm T_\pm$, and $R_0$
connects $\partial_0T_-$ to $\partial_0T_+$ while staying off $T_0$ (this
is the role of the push-off $E_0$). Put
\begin{equation}\label{eq:Omega}
\Omega=S_-\cup T_-\cup R_-\cup T_0\cup R_+\cup T_+\cup S_+.
\end{equation}
Ordered cyclically as in \eqref{eq:Omega}, \emph{consecutive} pieces meet
exactly along a common edge of length $\mu+\eps^2$ (the four shared edges
$\partial_- T_0,\partial_+T_0$ and the two copies of $E_0$). We check that
\emph{non-consecutive} pieces are disjoint for all small $\eps$. The three
triangles $T_-,T_0,T_+$ lie in the three distinct horizontal planes
$z=-\eps^3,\,0,\,+\eps^3$, hence are pairwise disjoint (their projections
to $z=0$ coincide, but their heights differ). Each connector has $z$-range
contained in $[-\eps^3,\eps^3]$ and is within horizontal distance
$O(\eps^3)$ of an edge of a triangle:
\begin{itemize}
\item $R_\pm$ joins $\partial_\pm T_0$ ($z=0$) to $\partial_\pm T_\pm$
($z=\pm\eps^3$); it is a $\eps^3$-thin rectangle hugging the
left/right edge, so it meets only its two neighbours $T_0,T_\pm$.
\item $R_0=S_-\cup S_+$ joins $\partial_0T_-$ and $\partial_0T_+$ to the
common push-off $E_0=\partial_0T_0+(0,\eps^3,0)$. Because $E_0$ is
displaced by $\eps^3$ in the $+y$ direction \emph{away from} $T_0$ (whose
interior lies in $y\le 0$ near that edge) and the rectangles $S_\pm$ have
short side $<2\eps^3$, the set $R_0$ stays in the slab
$\{y\ge -2\eps^3\}\cap\{|z|\le\eps^3\}$ and within horizontal distance
$<2\eps^3$ of the line $\partial_0T_0$; for $\eps$ small this region meets
$T_0$ only along $E_0$'s neighbours, i.e.\ $R_0\cap T_0=\emptyset$ except
that $R_0$ is glued to $T_\pm$ (its true neighbours) along
$\partial_0T_\pm$. In particular $R_0$ does not touch $T_0,R_\pm$.
\end{itemize}
Thus only consecutive pieces intersect, and only along the prescribed
edges. Hence $\Omega$ is an embedded piecewise-affine surface with
boundary.

Folding $T_-$ under and $T_+$ over $T_0$ realises the three planes; the
single seam $E_0$ is shared by $S_-$ (attached to $T_-$, below) and $S_+$
(attached to $T_+$, above). Tracing the cyclic boundary word
$S_-\,T_-\,R_-\,T_0\,R_+\,T_+\,S_+$ and back to $E_0$, the two arrows
along $E_0$ point in the \emph{same} direction (one inherited from below
via $T_-$, one from above via $T_+$), which is the orientation-reversing
gluing: $\Omega$ is an embedded M\"obius band, not a cylinder.

Finally triangulate each of the four connector rectangles by a single
diagonal. The resulting thin triangles have two long sides of length
$\ge\mu+\eps^2$ and one short side of length $<2\eps^3$.

\paragraph{The domain $\Omega'=\Sig/G_\beta$.}
Call a clean triangle \emph{wide} if it is isosceles with apex on the
opposite boundary line and its distinguished (base) side longer than
$\mu$; then the distinguished side is the strictly longest side. We will
build a polygonal block $P\subset\Sig$ that tiles $\Sig$ under a glide
reflection $G_\beta$.

\smallskip
\emph{The three wide triangles.}
Choose a wide clean triangle $T_0'$ in $\Sig$ whose distinguished edge
lies on the top boundary $\R\times\{1\}$ with length
\[
b:=\mu+\tfrac32\eps^2\in(\mu+\eps^2,\mu+2\eps^2),
\]
of height $1$ (apex on the bottom boundary). Its two crosscutting
(slanted) sides have the common length
\begin{equation}\label{eq:slant}
\sigma:=\sqrt{1+(b/2)^2}.
\end{equation}
We record the key estimate
\begin{equation}\label{eq:slantbound}
\mu<\sigma<\mu+\eps^2 .
\end{equation}
Indeed, at $\eps=0$ we have $b=\mu$ and
$\sigma=\sqrt{1+\mu^2/4}=\sqrt{1+\tfrac13}=\sqrt{4/3}=\mu$; for $\eps>0$
the value of $b$ increases, so $\sigma>\mu$ (left inequality). For the
right inequality, since $\sigma$ is increasing in $b$ it suffices to bound
$\sigma$ at the right endpoint $b=\mu+2\eps^2$:
$\sigma^2=1+\tfrac14(\mu+2\eps^2)^2=1+\tfrac{\mu^2}4+\mu\eps^2+\eps^4
=\mu^2+\mu\eps^2+\eps^4$ (using $1+\mu^2/4=\mu^2$), whereas
$(\mu+\eps^2)^2=\mu^2+2\mu\eps^2+\eps^4$; the difference is
$(\mu+\eps^2)^2-\sigma^2=\mu\eps^2>0$, so $\sigma<\mu+\eps^2$. This proves
\eqref{eq:slantbound} (and in particular $\sigma<\mu+\eps^2$ even at the
chosen $b=\mu+\tfrac32\eps^2$, by monotonicity).

Place two further wide clean triangles $T'_\pm$ congruent to $T_0'$ on
either side of $T_0'$; by the clean-triangulation convention their
distinguished edges lie on the bottom boundary $\R\times\{0\}$ (apices up).

\smallskip
\emph{The four parallelograms.}
Insert between $T'_0$ and $T'_\pm$ a parallelogram $R'_\pm$, and outside
$T'_\pm$ a parallelogram $S'_\pm$, all spanning the full height $1$ of the
strip, as follows. Each parallelogram has its two horizontal sides on the
two boundary lines $y=0,1$ and its two slanted sides shared with the
adjacent triangle edges. Consequently:
\begin{itemize}
\item the two slanted (long) sides of every parallelogram are crosscutting
edges shared with neighbouring wide triangles, hence have the common
length $\sigma$ of \eqref{eq:slant};
\item the two horizontal (short) sides of $R'_\pm$ have length $2\eps^2$,
and those of $S'_\pm$ have length $\eps^2$. These horizontal sides lie in
$\partial\Sig$ and will be the distinguished edges of the thin triangles
below.
\end{itemize}
This is geometrically realisable: a parallelogram of height $1$ whose
slanted sides have length $\sigma$ has horizontal offset
$\sqrt{\sigma^2-1}$ between its top and bottom edges (independent of the
short side length), and we are free to choose the horizontal short-side
length to be any positive number, in particular $\eps^2$ or $2\eps^2$.
Adjacent pieces are glued along full slanted edges of equal length
$\sigma$, so the union
\[
P=S'_-\cup T'_-\cup R'_-\cup T'_0\cup R'_+\cup T'_+\cup S'_+
\]
is a simply-connected polygonal region spanning the full height of the
strip, bounded on the left by the (left, slanted) outer edge $L_-$ of
$S'_-$ and on the right by the (right, slanted) outer edge $L_+$ of
$S'_+$. Both $L_-,L_+$ are slanted segments of length $\sigma$ joining the
two boundary lines, and they are parallel translates of one another (every
slanted edge in $P$ has the same direction vector, namely
$(\sqrt{\sigma^2-1},\,\pm1)$ with a consistent sign pattern dictated by
the alternating up/down apices; in particular $L_-$ and $L_+$ are
\emph{anti}-parallel as oriented boundary edges, which is what a glide
reflection requires).

\smallskip
\emph{The glide reflection and the fundamental-domain property.}
Let $w>0$ be the horizontal distance between the lines containing $L_-$ and
$L_+$ measured along $y=1/2$, i.e.\ the length of the intersection of $P$
with the centreline; we compute $w=\beta$ below. Define
\[
G_\beta(x,y)=(x+\beta,\,1-y).
\]
By construction $G_\beta$ is the unique orientation-reversing isometry of
$\Sig$ carrying the oriented left edge $L_-$ of $P$ onto the oriented right
edge $L_+$ of $P$: the reflection $y\mapsto1-y$ matches the slopes (which
are negatives of each other across the centreline, as just noted), and the
horizontal translation by $\beta$ carries $L_-$ to $L_+$. Hence
$G_\beta(L_-)=L_+$ and $P\cap G_\beta(P)=L_+$ is exactly the shared
slanted edge; likewise $P\cap G_\beta^{-1}(P)=L_-$. Because $P$ is a
height-$1$ strip block whose left and right boundaries $L_-,L_+$ are
identified by $G_\beta$ and whose interior is disjoint from all its
$G_\beta$-translates, the translates $\{G_\beta^k(P)\}_{k\in\Z}$ tile
$\Sig$ with pairwise-disjoint interiors and $\bigcup_kG_\beta^k(P)=\Sig$
(each translate is the next block to the right/left, glued along $L_\pm$).
Thus $P$ is a fundamental domain for $\langle G_\beta\rangle$ and
$\Omega'=\Sig/G_\beta=P/(L_-\!\sim\!L_+)$ is a flat M\"obius band, with the
seven pieces of $P$ inducing a clean triangulation of $\Omega'$ once the
parallelograms are cut (below).

\smallskip
\emph{The bound $\beta<\sqrt3+\eps$.}
The core length $\beta$ equals $w$, the length of the centreline chord
$P\cap(\R\times\{1/2\})$, which is the sum of the centreline chords of the
seven pieces. A wide triangle of base $b$ and height $1$ has, at half
height $y=1/2$, width $b/2\le(\mu+2\eps^2)/2=\tfrac\mu2+\eps^2$ (the width
interpolates linearly from $b$ at the base to $0$ at the apex). A
parallelogram with horizontal short side $s$ has centreline chord of
length exactly $s$ (its horizontal cross-section has constant width $s$).
Hence
\begin{equation}\label{eq:beta}
\beta=w= 3\cdot\tfrac{b}{2}+\big(2\cdot\eps^2+2\cdot 2\eps^2\big)
=3\Big(\tfrac\mu2+\tfrac34\eps^2\Big)+6\eps^2
=\tfrac{3\mu}{2}+\tfrac{33}{4}\eps^2=\sqrt3+\tfrac{33}{4}\eps^2<\sqrt3+\eps,
\end{equation}
using $b=\mu+\tfrac32\eps^2$ (so $\tfrac{b}{2}=\tfrac\mu2+\tfrac34\eps^2$),
$\tfrac32\mu=\sqrt3$, and the bound $\eps<1/100\Rightarrow
\tfrac{33}{4}\eps^2<\tfrac{33}{4}\cdot\tfrac{\eps}{100}<\eps$.
(The three wide triangles contribute centreline chord $\tfrac{b}{2}$ each,
the two $S'_\pm$ contribute $\eps^2$ each, and the two $R'_\pm$ contribute
$2\eps^2$ each; this is an exact equality, not merely an upper bound.)

\paragraph{The map $\overline f$: well-definedness, continuity, injectivity.}
Triangulate the four parallelograms of $P$ using the \emph{short}
diagonals. Concretely, write a parallelogram with bottom edge
$[\,(0,0),(s,0)\,]$, top edge $[\,(c,1),(c+s,1)\,]$ where
$c=\sqrt{\sigma^2-1}>0$ is the horizontal offset, and short side
$s\in\{\eps^2,2\eps^2\}$; the two diagonals have lengths
$\sqrt{(c+s)^2+1}$ and $\sqrt{(c-s)^2+1}$, and the short one is
$\delta:=\sqrt{(c-s)^2+1}$. Since $0<s<c$ (indeed $c=\sqrt{\sigma^2-1}$ and
$\sigma>\mu>1$ give $c>\sqrt{\mu^2-1}=\sqrt{1/3}>0.57\gg s$), we have
$0\le c-s<c<\sigma=\sqrt{c^2+1}$, hence
\begin{equation}\label{eq:diag}
\delta=\sqrt{(c-s)^2+1}<\sqrt{c^2+1}=\sigma<\mu+\eps^2 .
\end{equation}
Thus every short diagonal has length $\delta<\mu+\eps^2$, and the short
diagonal is strictly shorter than the slanted sides $\sigma$. Cutting each
parallelogram along its short diagonal produces two clean thin triangles,
each having: one horizontal short edge in $\partial\Sig$ (the
\emph{distinguished} edge, of length $s\in\{\eps^2,2\eps^2\}$), one slanted
crosscutting edge of length $\sigma$, and the diagonal as the second
crosscutting edge of length $\delta$.

This makes the triangulation of $\Omega'$ \emph{combinatorially identical}
to that of $\Omega$: both consist of three "big" triangles ($T'_k$ resp.\
$T_k$) and eight thin triangles (two per connector), glued in the same
cyclic incidence pattern, with the same vertex--edge--face adjacencies.
Fix the combinatorial isomorphism that sends $T'_k\mapsto T_k$ and each
thin domain triangle to the corresponding thin range triangle, matching
distinguished edge to distinguished (boundary) edge, slanted crosscut to
long ($\mu+\eps^2$) rectangle edge, and short diagonal to the
corresponding range diagonal.

Define $\overline f:\Omega'\to\Omega$ to be \emph{affine on each triangle},
sending the three vertices of each domain triangle to the three vertices of
the corresponding range triangle (in the matched order). This determines
$\overline f$ on each closed triangle uniquely. We verify the three
properties.

\smallskip
\emph{Well-defined and continuous.} The vertex assignment is consistent:
each vertex $v$ of the triangulation of $\Omega'$ is assigned a single
point of $\Omega$ (the corresponding vertex), independently of which
incident triangle we use, because the combinatorial isomorphism is a
bijection on vertices. If two triangles $\tau_1,\tau_2$ of $\Omega'$ share
an edge $e=[u,v]$, then the corresponding range triangles share the
corresponding edge $[\,\overline f(u),\overline f(v)\,]$, and the two
affine maps $\overline f|_{\tau_1},\overline f|_{\tau_2}$ agree on all of
$e$ because an affine map of a segment is determined by its values at the
two endpoints, and both restrictions send $u\mapsto\overline f(u)$,
$v\mapsto\overline f(v)$. Hence the affine pieces match along every shared
edge, so $\overline f$ is single-valued and continuous on
$\Omega'=\bigcup\tau$ (gluing lemma for finitely many closed sets meeting
along common faces).

\smallskip
\emph{Injectivity.} We use that $\Omega\subset\R^3$ is embedded (proved in
``The range $\Omega$'' above: distinct pieces are disjoint except along the
prescribed shared edges, so the only identifications in $\Omega$ are those
of the abstract gluing). The map $\overline f$ is the composition of (a)
the abstract piecewise-affine homeomorphism from $\Omega'$ to the abstract
M\"obius band $M$ obtained by gluing the eleven model triangles along
matched edges, with (b) the realisation map $M\to\Omega\subset\R^3$. Map
(a) is a homeomorphism by construction (it is affine and bijective on each
triangle and respects the gluing). Map (b) is injective precisely because
$\Omega$ is embedded: the only points of distinct model triangles that get
identified in $\R^3$ are those lying on shared edges, which are already
identified in $M$. Hence $\overline f$ is injective. (Concretely: if
$\overline f(p)=\overline f(q)$ with $p\in\tau_1,q\in\tau_2$, then the
common image point lies in $T_{i}\cap T_{j}$ for the corresponding range
triangles; by embeddedness this intersection is either empty, a shared
vertex, or a shared edge, and on each affine piece $\overline f$ is a
bijection onto its image triangle, forcing $p=q$.)

\paragraph{Verification of (i)--(iii) and the squeeze property.}
Let $f=\overline f\circ\pi:\Sig\to\Omega\subset\R^3$, where
$\pi:\Sig\to\Omega'$ is the quotient map. Then:
\begin{itemize}
\item \textbf{(i) Invariance.} By construction $f=\overline f\circ\pi$
and $\pi\circ G_\beta=\pi$, so $f\circ G_\beta=f$.
\item \textbf{(ii) Injectivity on orbits.} $\overline f$ is injective and
$\pi(p)=\pi(q)\iff q\in G_\beta\cdot p$ (since $P$ is a fundamental domain
and the only boundary identification is $L_-\sim L_+$ via $G_\beta$), so
\[
f(p)=f(q)\iff\overline f(\pi p)=\overline f(\pi q)\iff\pi p=\pi q\iff
q\in G_\beta\cdot p .
\]
\item \textbf{(iii) Squeezing.} We verify the \emph{strict} squeeze
property (Convention \ref{conv:squeeze}) on every triangle: the
distinguished (boundary) edge is strictly shortened and the two
crosscutting edges are strictly lengthened. In particular this gives
condition (iii) (strict shortening of boundary edges).
\begin{itemize}
\item \textbf{Wide triangles $T'_k$ ($k\in\{0,+,-\}$).} The distinguished
(base) edge has length $b=\mu+\tfrac32\eps^2>\mu+\eps^2$, and the two
crosscutting (slanted) edges have length $\sigma<\mu+\eps^2$ by
\eqref{eq:slantbound}. The map $\overline f$ sends all three edges to the
three edges of the equilateral triangle $T_k$, each of length exactly
$\mu+\eps^2$. Hence:
\[
\text{distinguished: } b=\mu+\tfrac32\eps^2\ \mapsto\ \mu+\eps^2
\quad(\text{shortened, since }\tfrac32\eps^2>\eps^2),
\]
\[
\text{crosscutting: } \sigma\ \mapsto\ \mu+\eps^2
\quad(\text{lengthened, since }\sigma<\mu+\eps^2\text{ by \eqref{eq:slantbound}}).
\]
\item \textbf{Thin connector triangles.} Such a triangle has
distinguished (boundary) edge of length $s\in\{\eps^2,2\eps^2\}$ (so
$s\ge\eps^2$), and crosscutting edges of lengths $\sigma$ (slanted) and
$\delta$ (short diagonal), \emph{both} $<\mu+\eps^2$ by
\eqref{eq:slantbound} and \eqref{eq:diag}. The corresponding range thin
triangle (from cutting a rectangle of long side $\mu+\eps^2$ and short
side $\le2\eps^3$) has distinguished edge of length $\le2\eps^3$, slanted
crosscut of length $\mu+\eps^2$, and diagonal crosscut of length
$\sqrt{(\mu+\eps^2)^2+(2\eps^3)^2}\ge\mu+\eps^2$. Hence:
\[
\text{distinguished: } s\ge\eps^2\ \mapsto\ \le2\eps^3
\quad(\text{shortened, since }2\eps^3<\eps^2\text{ for }\eps<\tfrac12),
\]
\[
\text{slanted crosscut: } \sigma<\mu+\eps^2\ \mapsto\ \mu+\eps^2
\quad(\text{lengthened}),
\]
\[
\text{diagonal crosscut: } \delta<\mu+\eps^2\ \mapsto\ \ge\mu+\eps^2
\quad(\text{lengthened}).
\]
\end{itemize}
In every triangle the distinguished/boundary edge is strictly shortened
and both crosscutting edges are strictly lengthened, so $f$ is a strict
squeeze map. (The two strict-shortening estimates $b>\mu+\eps^2$ and
$2\eps^3<\eps^2$, and the three strict-lengthening estimates
$\sigma<\mu+\eps^2$, $\delta<\mu+\eps^2$, are all verified above for
$\eps\in(0,1/100)$.)
\end{itemize}
Thus $\beta<\sqrt3+\eps$ is realized, proving Statement 2. \hfill$\square$

%=====================================================================
\section{Conclusion}

By Statement 1 every realized $\beta$ satisfies $\beta>\sqrt3$, so
$\inf\{\beta>0:\beta\text{ realized}\}\ge\sqrt3$. By Statement 2, for
every $\eps>0$ some realized $\beta$ satisfies $\beta<\sqrt3+\eps$, so the
infimum is $\le\sqrt3$. Hence
\[
\boxed{\ \inf\{\beta>0:\ \beta\text{ is realized}\}=\sqrt3\ }
\]
and the infimum is not attained.

\end{document}
